% 安全性分析的补充细节。
\section{原始执行关联的细节}
\label{supp:original-execution}

\noindent\textbf{最早的冲突 polka。}
固定一个高度，以及第 $r$ 轮中针对非空值 $v$ 的一份提交证书，并从该证书中选出两名诚实签署者 $K$。每名签署者在发出其 precommit 之前先锁定 $v$。证书可以在这些签署者已经推进之后才汇集，所以它被观察到的时间无关紧要。

假设存在针对不同值的更高轮 polka（若 nil polka 允许解锁，则也计入）。取存在这种 polka 的最小轮次 $s>r$，并在 $s$ 内取其不同的 prevote 全部存在的最早事件 $\tau$。其三名投票者中包含某个 $m\in K$。诚实节点的轮次递增，所以 $m$ 在第 $s$ 轮的 prevote 晚于它在第 $r$ 轮的 precommit，而没有解锁理由时，$m$ 必须对其已锁定的值做 prevote。在中间轮次中重新锁定 $v$，并不改变这一约束。来自 $s$ 以下某一轮的冲突理由，与 $s$ 的选取相矛盾；来自第 $s$ 轮的冲突理由，必然在 $m$ 投出 $\tau$ 时被计入的那一票之前就已存在，这与 $\tau$ 的选取相矛盾。来自 $s$ 之上的证据，不可能先于第 $s$ 轮中较晚的投票。因此这样的 polka 不存在，而更高轮的冲突提交必然需要它。同一轮中冲突的提交，则需要某个诚实节点做出两次 precommit。这一论证依赖于所述的锁定转移规则和轮次单调性规则，并不要求及时收到提交证据。

\noindent\textbf{配对、验证与执行。}
对候选更新做归纳。本地构造把区块体与其原始分片配对，远程补全则由经检查的分片重建区块体。在 Locked 或 Valid 值移动时，二者一并转移；切换目标则清除旧区块体。历史安装不触及这些活跃候选。非空的 precommit 紧随匹配的 polka。诚实的 prevote 在实际的提案验证或更早的锁定中有其理由，追溯后者时锁定轮次逐步降低，直到抵达验证。最后的完整身份守卫在 Finalize 之前固定原始交易、高度、时间和其他应用上下文。同步与重放提供相同的原始材料，所以在这一调用边界之内，同一高度的重试不会通过稳定哈希缓存引入第二个原始区块体。

\noindent\textbf{写集。}
正文中的证明依赖于下列写集。赔付决定写入：签发组织备付金账户、义务状态、Promise 与覆盖记录（Paid 和 Remaining）、授权额度用量（Spent）、公共缺口、消费方的支付记录（Pending 和费用阶段）、奖励与退款余额及费用输出、到期索引条目的移除，以及决定记录及其决定索引条目和待历史表示条目。赔付之后源交易的执行写入偿还：签发组织备付金账户、义务状态（Recovered）、覆盖记录（Recovered）和授权额度用量（Spent）。它不改变决定记录。补交任务只写入成员本地的 outbox。历史表示命令或批处理写入逻辑修订、头命令、任务或批处理任务及其逐输出引用、队列条目，以及待历史表示条目的移除。安装只写入本地区块存储的版本及其自身的进度游标。历史表示和安装的写入都不属于经济投影。

\noindent\textbf{固定目标集合的收敛性。}
固定一个原始区块体和一组已决定槽位。每个槽位的替换消息是网络和输出的函数，因此无论如何打包，都以相同的消息替换相同的槽位。若变色龙关系对每个固定的消息与承诺都有唯一的有效 opening（参数正确生成且公开指数与群的阶互素时即成立），则每个槽位 opening 和每个分片 opening 都是其关系的唯一解。分片上下文取决于高度和分片索引，与修订计数器无关。因此，只要覆盖该集合，历史表示命令的任何分组或顺序都得到相同的区块体和分片 opening。单条目命令、单个批处理以及中间规模的批处理，在命令标识、修订号和应用哈希上各不相同。该陈述针对的是最终字节，并如上所述依赖于参数生成和编码。

\noindent\textbf{跟随者更新。}
对于从相同的本地状态、配置和游标出发的跟随者，只要两种表示都能解码，相同的固定支付字段与经认证的结果就诱导出相同的经济更新。可变的资金引用与 opening 不选择任何 Paid、Pending 或费用效果。决定被完整认证，执行结果承诺则认证其效果和费用输出。因此，经济观察与对最新可变历史区块体的独立验证保持分离。

\noindent\textbf{分析中使用的关系。}
三种关系比较的是不同的对象。历史表示的经济惰性，比较一次历史表示或安装步骤前后的经济投影，并投影掉命令标识、修订计数器和物化任务。固定目标收敛性，比较不同的历史表示命令分组与顺序为同一原始区块体和同一已决定集合所产生的最终区块体和分片 opening。重放相等性固定整个原始命令序列及其区块上下文，确定性执行随后重现每个应用哈希。跨历史比较固定经济命令、公共时间和历史目标坐标，因此投影相等并不能确定整个应用状态相同，也不允许删除区块或给区块重新编号。这三种关系都不比较以不同顺序安排父交易到达与赔付的历史；正文的延迟中性定理在其自身的假设下做这一比较，它等同的是最终的 CAL 持有者和备付金余额，而不是应用哈希或修订元数据。

\noindent\textbf{本地失效的写集与重放。}
失效判断仅发生在应用认证区块中执行成功且确实生效的付款记录时；
\texttt{INSTALL}不能触发失效。在Follower覆盖Spend记录之前，
它读取原Candidate，并核对不可变Approval以及精确的本金输入（最终或认证输入）或
用户费用输入实例。若旧事实已经被观察、已有安装证书材料或outbox，
或已经发生结算、Pending、赔付、回款或费用效果，则将其视为会计
状态矛盾并停止处理，而不是释放容量。

新增的失效写集仅包括受影响的本地资源Slice和Approval的LocalProgress。
对每项Debit，它从Reserved减去
$\mathsf{Cap}-\mathsf{Applied}$，向Available增加同量，将累计
Applied置为Cap，并记录Invalidated、首个触发失效的公共事实及高度。
不可变Approval和其他输入锁保持不变。同一Overlay使同一区块中后续
冲突输入或交易不再重复回收；Follower游标抑制区块重放，累计
Applied向量防止重启后再次入账。成员存储Schema相应升级，避免旧
二进制忽略新的失效终态。这些变更与成功消费记录和跟块游标一次提交。Schema 8拒绝旧成员库；实验使用新目录，不进行迁移。

独立输入上已经形成QC、但因全局Intent绑定而永久被拒绝的来源不适用
上述失效规则：其QC责任、输入锁和签发容量继续保留。当前Relay可能
持续重试该来源；终止此类重试属于独立的工程后续工作，并不得因此
释放已获证责任。

\subsection{实际入口与费用证据}
\label{supp:interface-fee-evidence}
工件中的 2026-10-04 接口审计保存源码指纹、命令与断言日志。测试补充基于检出版本 \texttt{863952d}；节点生产逻辑相对被测基线 \texttt{f856c7b} 不变。本地回归使用 Windows amd64 上的 Go 1.27.1，提供正确性证据，不替换性能测量。

\noindent\textbf{涉及身份的转移。}
固定版本的 CometBFT overlay 将八处区块体哈希比较替换为 \texttt{matches\allowbreak BlockID}，同时检查体哈希与配对的部件集头。剩余 Finalize 哈希检查位于部件头检查之后，投票扩展调用方已受完整身份匹配守卫约束。\texttt{Test\allowbreak Stable\allowbreak Hash\allowbreak Identity} 调用真实共识状态方法，覆盖提交目标替换、旧部件、精确最终化、Locked/Valid 配对、POL 轮次，以及法定票先于或后于区块体到达。polka 先于部件的场景还从 Propose 步骤开始：补齐候选后只针对精确身份预提交一次，此后的尝试既不产生迟到预投票，也不产生第二次预提交。独立的 nil 预提交场景确认迟到部件不产生第二次投票。这同时检查新增的当前轮 polka 转移与比较守卫。

\noindent\textbf{原始材料与应用认证。}
\texttt{Test\allowbreak Consensus\allowbreak Rejects\allowbreak Relabelled\allowbreak Opening} 将重标记为修订零的兼容适配送入共识消息的 \texttt{Validate\allowbreak Basic} 与 \texttt{Msg\allowbreak From\allowbreak Proto}；原始准入拒绝非原始开口。部件集的包含性验证本身会接收已签名根下的重标记部件；对非原始开口的拒绝发生在共识消息边界。真实 BlockStore 改写/重放与追块测试确认修订后仍保留、提供原始区块体和部件。\texttt{Test\allowbreak Public\allowbreak Authorization\allowbreak At\allowbreak ABCIEntry} 先预热有效验证缓存，再检查 \texttt{CheckTx}、\texttt{ProcessProposal} 与 \texttt{FinalizeBlock} 拒绝且业务状态不变。伪 QC 场景中，解码、所有者授权、准入准备和摘要绑定均通过，随后 QC 认证返回 \texttt{ErrAuth}；无关 QC 与伪所有者场景在更早的守卫拒绝。独立的缓存测试要求复用字节完全相同，并重新检查当前账本有效性。

\noindent\textbf{托管与关闭。}
\texttt{PrepareDirectVector} 以溢出检查核算基础费用及每个认证输入的一次赔付费用。\texttt{feeStage} 应用累计托管差量，\texttt{endObligation} 只在离开 Open 时允许一次赔付扣费，\texttt{closeFee} 在准入上限内核算实付费用。\texttt{Test\allowbreak Fee\allowbreak Closure\allowbreak Multiple\allowbreak Obligations} 包含十二个规则层场景：两种付费模式、零至两个赔付输入、最低或富余托管。断言覆盖每步托管与资产守恒、精确最终费用和退款、Pending 归零时关闭及重复操作幂等。\texttt{Test\allowbreak Fee\allowbreak Closure\allowbreak Fulfillment\allowbreak Before\allowbreak Compensation} 的另外两个场景使用不同父交易：一个正常到达，另一个随后赔付触发消费方关闭，第二个父交易迟到不改变已关闭托管。公开决定与多槽位测试另行覆盖应用边界。成员跟块测试在赔付与来源到达后比较本地托管与公共支付状态，并检查累计恢复和已消耗 FUEL 不循环恢复。

\noindent\textbf{观测所支持的结论。}
这些测试将证明中的转移假设对应到实际入口，不替代锁定或会计论证。无义务决定守卫拒绝且不返回变更或效果，对应服务契约不提供仅凭证书兑付。最终混合负载资本审计读取每轮八名成员的已记录计数和资源终点：三轮资源不足计数均为零，最终 CAL/工作资源预留为零，而实际 FUEL 支出仍保留。该终点审计与未公开证书的理论覆盖、持续准入的负载条件分别陈述。
